\section{Reproducible protocol specification}\label{sec:protocol-spec}
\subsection{Security experiments and output leakage}\label{sec:security-experiments}
Let an approved instance $I$ contain the catalog root, integer model,
quantization, reference centers, circuit shapes and verification-key hashes.
The owner retains $I$ independently of the prover. A query authorization
contains the complete candidate centers and catalog scope, rather than a
prover-supplied label identifying them. The verifier's accepted decision is
computed from checked public endpoints, not a supplied decision string.

\begin{definition}[False-decision experiment]
An adversarial prover commits the original and derived catalogs and moments.
An independent challenger samples a uniform nonzero field challenge after
that commitment and pins it to $I$. The prover may submit query transcripts
for externally approved center sets and any claimed refinements. It wins if
the public verifier accepts an interval excluding the true integer cost
change, or a strict sign different from that cost change, on the full fixed
catalog. Missing, duplicate or substituted coverage is rejection.
\end{definition}
Conditioned on binding commitments and argument soundness, the proposition's
proof bounds this event by the fingerprint error for the fixed session.
The row relation fixes original coordinates and minimizing old groups; the
mask-chain relation makes the folded totals equal the committed moment
polynomials. A false integer moment vector has a nonzero discrepancy
polynomial because the checked ranges exclude field aliases. Outside a
polynomial collision, exact assigned change and conservative row bounds
imply interval containment. Refinement preserves containment by same-index
openings, canonical endpoint reconstruction and unique indices. An attacker
who changes an authorized context must instead break a binding or proof
assumption. The argument covers a bounded sequential session under these
assumptions; it is not a universally composable protocol theorem.

\begin{definition}[Disclosed transcript and leakage-relative privacy]
For instance $I$ and authorized queries $C^{1: m}$, let $\mathcal L$ contain
all public metadata, commitment roots, challenge, public masked row
contributions and mask hashes, assigned aggregates, batch intervals,
refinement indices and corrections, and any fallback aggregate costs.
It also contains the authorized queries, termination path and proof counts.
In the real experiment an auditor receives these values and honest proofs;
in the simulated experiment it receives the same values and simulated
proofs for their corresponding public statements. Leakage-relative privacy
means that no efficient auditor distinguishes these experiments beyond the
underlying proof system's zero-knowledge advantage.
\end{definition}
This is a conditional proof-transcript property. Sequential simulation
replaces each proof for its fixed public statement; it neither simulates
nor hides $\mathcal L$ from the auditor. Public masked contributions are
included explicitly, so no unproved masking-composition claim is needed.
For adaptive authorized queries the leakage grows with every answer.
Equation~\eqref{eq:probe-leakage} proves that group-moment hiding would fail
under that stronger experiment. The owner must assess this disclosure before
authorizing queries. A query budget alone would bound the number of answers
but would not prevent recovery from a sufficiently revealing allowed set.
Network timing, host side channels and a malicious owner intentionally
publishing records lie outside this proof-transcript experiment.

\subsection{Model, commitments and challenge lifecycle}
The instance fixes $N,D,k$, the quantizer and clipping range, integer basis
$B$, scale $T$, anchor $a=c^0_0$ and deployed centers $C^0$. The prototype
supports power-of-two $N\leq2^{20}$, $D\leq10,000$, $2\leq k\leq128$,
rank at most 128, input bit bound at most 16, and $1\leq T\leq2^{20}$;
these are implementation guards rather than evaluated scales.
All commitments use the instantiated Poseidon2 hash with length fixed by
the model schema. Original-row leaves use domain 101, $N,D$, tag zero,
index, salt and coordinates. Derived leaves use domain 102 and additionally
bind the model hash, with values $(H_i,E_i,g_i)$. Tree nodes use domain 105
and level; domain 106 binds the model hash and catalog roots into a common
context. Reserved energy and tail roots are zero for the selected profile.
The moment commitment uses domain 131, model hash, original/derived/group
roots, salt, and the complete ordered counts and coordinate sums.

The data owner retains the approved model, original root and commitment
receipt independently of the prover's public bundle. After the original,
derived/membership and moment commitments are fixed, the auditor samples
$r$ uniformly from $\mathbb F_p\setminus\{0\}$ and retains its receipt.
Here $p$ is the BN254 scalar-field modulus. The benchmark uses the operating
system's cryptographic randomness and explicit local role separation.
Approved verification-key hashes come from the setup authority for the
specified model and circuit shapes, not from untrusted prover metadata.
Every accepted proof uses the same pinned roots, model, challenge and keys.
The challenge is reused within this one authenticated preprocessing session.
Changing commitments or redoing preprocessing creates another session and
requires another independently issued challenge and accounted work.

The verifier rejects zero or noncanonical field values, substituted receipts,
incomplete or duplicated catalog coverage, reordered indices, broken mask
chains, missing query batches and centers different from the externally
approved task. The prototype does not implement authenticated network
transport, a public randomness service or an unbounded retry/grinding model.
These are explicit assumptions of the conditional soundness statement.

\subsection{Indexed mask chain and moment equality}
For each $i\in\{0,\ldots,N\}$ and group $g$, sample independent private
$\mu_{ig}\in\mathbb F_p$. Define
\begin{align}
 P_i(r)&=1+\sum_{d=0}^{D-1}x_{id}r^{d+1},\\
 z_{ig}&=\mathbf1[g_i=g]P_i(r)+\mu_{i+1,g}-\mu_{ig}.
\end{align}
The public mask hash at position $i$ binds domain 132, model hash, original,
derived, group and moment roots, $r$, $i$ and the complete mask vector
$\mu_i$. Row $i$ proves its original and derived openings, exact old argmin,
both adjacent mask hashes and every $z_{ig}$. The public verifier requires
indices $0,\ldots,N-1$ exactly once and equality of adjacent hashes.
The moment-link relation opens the count/sum commitment and endpoint masks
and checks, for every group,
\begin{equation}
 \sum_i z_{ig}=n_g+\sum_{d=0}^{D-1}S_{gd}r^{d+1}
                       +\mu_{Ng}-\mu_{0g}\pmod p.
\end{equation}
Thus internal masks telescope. Their being chosen after $r$ does not remove
the moment-equality condition: endpoints are linked to the same authenticated
chain and checked in the moment proof. Counts obey $0\leq n_g\leq N$,
$\sum_g n_g=N$, and sums obey $|S_{gd}|\leq NX_{\max}$ for coordinate bound $X_{\max}$.
A fixed false moment vector gives a nonzero discrepancy polynomial of degree
at most $D$. Binding before the independent challenge permits the polynomial
root argument in Section~\ref{sec:security}. The single-session event is
independent of empirical prediction agreement.

\subsection{Integer limits and signed arithmetic}
Let $X_{\max}=2^{b}$ bound every input and center coordinate,
$W_0=2X_{\max}$, rank be $s$, and $\lambda_j=\sum_d|B_{jd}|$.
With $|B_{jd}|\leq T$, define
\begin{align}
 H_{\max}&=TW_0\max_j \lambda_j,\\
 R_{\max}&=W_0\max_d\left(T^2+\sum_j|B_{jd}|\lambda_j\right),\\
 E_{\max}&=D R_{\max}^2,\qquad \rho=\lceil\sqrt{E_{\max}}\rceil.
\end{align}
These bounds apply to both rows and new centers relative to the bounded
anchor. A row-radius witness obeys $E_i\leq\rho_i^2$ and $0\leq\rho_i\leq\rho$;
native construction uses the smallest ceiling radius. Pairwise center
residual radii are at most $2\rho$. For $\gamma=\max_a\sum_b|G_{ab}|$, the
implementation bounds public affine coefficients and constants by
\begin{align}
 C_{\max}&=4H_{\max}(\gamma+T^2),\\
 K_{\max}&=T^2(2sH_{\max}^2+2E_{\max})
                         +2\gamma sH_{\max}^2,\\
 \Lambda&=sH_{\max}C_{\max}+K_{\max}+4T^2\rho^2+1.
\end{align}
Here $Q$ is the maximum row count in a switching batch; the selected profile
uses $Q=16$. The switching comparison guard uses $L_G=16\Lambda$, each signed batch endpoint
is bounded by $2Q L_G$, and compilation requires
$64(N+Q+k)L_G<p$. The model guard separately checks its head, energy,
score and distance limits multiplied by $16N$ against $p$. Exact old
linear scores have magnitude at most $3DX_{\max}^2$; their differences are bounded
by $6DX_{\max}^2$. The assigned-change range is $|A|\leq4NDX_{\max}^2$.

Signed values are range-checked after an explicit nonnegative shift, with
binary/range decomposition and bounded integer comparisons; field residues
alone are not treated as signed integers. Public centers are canonical
bounded integers. Roots, challenge and masked contributions are canonical
field elements. The public verifier derives affine parameters independently
and accumulates $T^6A$ and batch endpoints with arbitrary-precision integers.
The machine-readable export retains the exact model parameters and limits;
an out-of-guard model is excluded before accepted proof generation rather
than silently reduced modulo $p$.

\subsection{Refinement relation and verifier obligations}
The refinement relation uses one row per proof. It opens both approved
original and derived leaves at the public row index. For $E_i=0$ it forces
$\bar\rho_i=0$; otherwise it enforces
$(\bar\rho_i-1)^2<E_i\leq\bar\rho_i^2$. It reconstructs every canonical
switching comparison using the approved center coefficients and center
residual radii. Original-space linear scores
$\|c_j\|^2-2\langle x_i,c_j\rangle$ determine the exact new minimum and
the old group's new score, hence $b_i$. Their checked magnitudes and
differences obey the $3DX_{\max}^2$ and $6DX_{\max}^2$ limits above. Canonical endpoints
and $T^6b_i$ use the same signed switching bounds as the base relation;
each correction has magnitude at most $4QL_G$ and its required sign is
checked rather than interpreted from bare
field residues. Minimal radii and canonical parameters are essential to
the monotone endpoint replacement argument.

An independently pinned refinement key authenticates this extra relation;
it cannot replace or overwrite the base key family. The verifier replays
all base proofs under their original pin before accepting any correction,
reconstructs public coefficients from the same approved centers, rejects
duplicated or out-of-range row indices and more than 16 corrections per
candidate, and recomputes the resulting interval and decision. Targeted real
proof checks reject altered corrections, duplicate indices, missing base
coverage, changed centers, substituted paths and modified proof bytes.
These checks complement the conditional argument; they are not a general
security audit.

\subsection{Supplementary advance-bank affine control}\label{sec:known-family-results}
This control allows the affine baseline to see the union of all proposed
center sets before authenticating its exact scores. Centers change between
queries, but all belong to this advance bank. It therefore cannot supply the
timing of a previously unseen-center service. Table~\ref{tab:known-family}
retains its measured results without placing them in the main timing ranking.
An online rebuild would have to charge new preparation, authenticated score
derivation, required keys, query proofs and replay when an unseen center
arrives. Completed request-local online rebuild controls are reported in
Table~\ref{tab:optimized-exact}; they are distinct measurements, not multiples
or relabelings of this advance-bank control.
\begin{table*}[tbp]
\centering
\caption{Supplementary exact affine measurements with advance knowledge of every future center. The bank includes all evaluated center sets before preprocessing; sets include the deployed reference. These are a different information contract from the main unknown-future-center task. No online-affine speed comparison is inferred. Assignments and costs are exact within the bank.}
\label{tab:known-family}
\begin{adjustbox}{max width=.98\textwidth,center=\textwidth}
\begin{tabular}{lrrrrrr}
\toprule
Cohort & $N$ & Workers & Bank & Sets & Min. & Proofs\\
\midrule
MNIST,19 & 128 & 2 & 85 & 20 & 40.174 & 288\\
MNIST mature & 128 & 8 & 15 & 5 & 4.710 & 168\\
Fashion early & 128 & 8 & 15 & 5 & 4.443 & 168\\
MNIST scale & 1024 & 8 & 25 & 5 & 44.704 & 1344\\
\bottomrule
\end{tabular}
\end{adjustbox}
\end{table*}


\subsection{Basis construction}
Let $r_0$ be the rank of the deployed-center differences. The centroid-informed
basis begins with an orthonormal SVD basis of
$\operatorname{span}\{c^0_j-c^0_0:j>0\}$. Subtract the anchor $c^0_0$ from
training rows, remove this span and mean-center the remaining training
residuals. Randomized PCA supplies $63-r_0$ directions, with two
reprojections against the centroid span and a nonpivoted QR preserving PCA
order. Column signs are canonicalized and the persisted basis is rounded at
scale $2^{14}$. The PCA instantiation fits all 63 prediction directions on the
same training split without the centroid-span prefix. The auxiliary
certificate appends the first 16 directions of a training-residual randomized
SVD with 64 components, four iterations and seed 20261025, rounded at the same
scale. No held-out audit row or future candidate fits either representation.

\subsection{Supporting completed controls and cost allocations}
The following tables retain the earlier two-worker controls, observed Full
timing variation and modeled cumulative allocations. They supplement the
main matched eight-worker comparison; ratios are never formed across worker
counts or hosts. Intermediate allocation and reset curves reuse recorded
components and are not fresh deployment measurements.
\begin{table*}[tbp]
\centering
\caption{Completed real-proof controls, $N=128$, two Go workers and one repetition per workload. Active minutes include cold preparation, setup and fresh public replay. Neither cohort triggers fallback. Future candidate centers are not required by preprocessing.}
\label{tab:completed-proofs}
\begin{adjustbox}{max width=.98\textwidth,center=\textwidth}
\begin{tabular}{llrrr}
\toprule
Cohort, updates & Method & Minutes & Full/method & Proofs\\
\midrule
MNIST, 19 & Exact Full & 196.703 & 1.000 & 320\\
MNIST, 19 & Centroid & 61.637 & 3.191 & 299\\
Fashion, 2 & Exact Full & 39.457 & 1.000 & 48\\
Fashion, 2 & Centroid & 24.771 & 1.593 & 146\\
\bottomrule
\end{tabular}
\end{adjustbox}
\end{table*}

\begin{table*}[tbp]
\centering
\caption{Centroid-basis certificate controls (rank 79; prediction rank 63): conclusive cost-change decisions, native predictor agreement (mean/worst), native assignment-certified fraction, native mean absolute relative cost error, and accepted change-interval width. Percentage columns concern different quantities; prediction rank remains 63.}
\label{tab:completed-quality}
\begin{adjustbox}{max width=.98\textwidth,center=\textwidth}
\begin{tabular}{lrrrrrr}
\toprule
Cohort & Decisions & Agr. & Worst & Label cert. & Cost error & Width\\
\midrule
MNIST & 19/19 & 99.794 & 99.219 & 94.901 & 0.0240 & 1.0250\\
Fashion & 2/2 & 100.000 & 100.000 & 100.000 & 0.0018 & 0.0000\\
\bottomrule
\end{tabular}
\end{adjustbox}
\end{table*}

\begin{table*}[tbp]
\centering
\caption{Sensitivity to observed Full timing variation on byte-identical integer task inputs. The second cold Full execution is the paid moments-only fallback, including fresh proofs and public replay; it is not an additional timing campaign. The ranges divide both observed Full totals by each single selected-method measured component sum. They are not confidence intervals or independently repeated paired comparisons.}
\label{tab:full-repeat}
\begin{adjustbox}{max width=.98\textwidth,center=\textwidth}
\begin{tabular}{lrrrr}
\toprule
Cohort & First Full min. & Repeat min. & Centroid ratio range & PCA ratio range\\
\midrule
MNIST mature & 20.971 & 18.758 & 1.745--1.950 & 2.006--2.242\\
\bottomrule
\end{tabular}
\end{adjustbox}
\end{table*}


\subsection{Data preparation and evidence boundaries}
MNIST and Fashion-MNIST~\cite{lecun1998,fashion2017} use pinned 5,000-row
source samples with sample seed zero, removal of zero rows, and a 60\%
training split with permutation seed $741+s$ for fit seed $s$. Input scaling
is fitted on training rows only, mapping its largest absolute coordinate to
4,095 and clipping quantized audit coordinates to $[-4096,4096]$.
The first 1,024 held-out records form the largest native catalog; smaller
catalogs are fixed prefixes. Training-only Lloyd snapshots define deployed
ages 1 and 8. Candidate snapshots are 2/4/8/16 or 9/12/16/24 as applicable.
All distinct supplied candidates are retained, including unfavorable,
inconclusive and below-target outcomes. The auxiliary directions are fitted
on training residuals; no audit outcome selects individual updates.
The original 19-update MNIST control uses seed 11 and all training Lloyd
snapshots through iteration 24, deduplicated to 19 nonidentity proposals;
the smaller fresh-fit candidate grids do not replace that original trace.

Proof controls run on an AMD Ryzen 7 5800H Windows laptop with eight physical
cores, 16 logical processors and approximately 13.9\,GiB usable RAM.
The completed original cohorts use two Go execution threads. New matched
challenge controls pin eight threads for every method and fallback, a
1,200\,MiB Go soft memory limit, a 3.5\,GiB process RSS cap and a 650\,MiB
free-memory reserve for the initial Full run. Subsequent phases use a
recorded scheduling amendment: Windows low-memory notification, at least
1\,GiB available commit, and launch headroom with retry hysteresis; the
same Go soft limit, worker count and RSS cap remain. Available RAM and
interrupted work differ across phases. These are descriptive component
timings, not a controlled memory-scaling experiment. Compact batches have 16 records and Full batches eight,
under a 1.6-million-constraint guard. All failed attempts and resume costs
remain in the ledger. Component microbenchmarks of parallelism or binary I/O
are not substituted for complete comparative deployment timings.

Public replay uses approved pins, centers, keys, proofs and public metadata;
it does not load original rows, private groups, moments or diagnostic labels.
Hashes check artifact integrity and targeted tests check specified
substitutions; neither is a complete security audit. The research artifact
remains private for supervisor review. The older server exports have been
recovered and transfer hashes checked, but their comprehensive fresh local
replay remains pending; they do not enter the current timing tables. The
additional current-cohort public-only replay has completed: 7,679 proofs
passed fresh verification in new processes, covering the matched controls,
supplementary advance-bank controls and the extra timing execution. Source
manifests and proof bytes were checked before and after replay. This used the
same frozen verifier implementations, one worker and a 256\,MiB Go soft memory
target; it is additional implementation checking, not an independently
implemented verifier or a new timing measurement. Its receipt is
\texttt{public\_replay\_round3\_20261026/status.json}. Per-control public replay
receipts are retained separately.
Earlier method variants are retained
in the development archive, outside this manuscript's comparative method set.

\subsection{Extended native sweep and implementation measurements}
The following numerical controls use the frozen profile and new source
splits described in the main evaluation. The age-1 scale series is primary;
the additional age-8 prefixes evaluated by the implementation are retained
as paired diagnostics. Neither additional prefixes nor the two bases add
independent splits. No catalog is enlarged by duplicating records.
\begin{table*}[tbp]
\centering
\caption{Native center-count sweep. Four seeds contribute to $k=5$; $k=3,10,20$ use seed 173. All declared distinct candidates are retained. Width is percent of deployed cost after bounded refinement; Below 99 counts requests missing the working agreement target. A 0/0 regression entry means no natural regression was present.}
\label{tab:optimized-centers}
\begin{adjustbox}{max width=.98\textwidth,center=\textwidth}
\begin{tabular}{lrlrrrrrr}
\toprule
Data & $k$ & Basis & Decisions & Regressions & Worst agr. & Label cert. & Width & Below 99\\
\midrule
Fashion & 3 & Centroid & 8/8 & 0/0 & 99.902 & 99.426 & 0.376 & 0\\
Fashion & 3 & PCA & 8/8 & 0/0 & 100.000 & 99.524 & 0.297 & 0\\
Fashion & 5 & Centroid & 28/28 & 5/5 & 99.805 & 98.982 & 0.227 & 0\\
Fashion & 5 & PCA & 28/28 & 5/5 & 99.707 & 98.525 & 0.212 & 0\\
Fashion & 10 & Centroid & 6/8 & 2/3 & 99.805 & 97.925 & 0.412 & 0\\
Fashion & 10 & PCA & 6/8 & 2/3 & 99.805 & 96.680 & 0.413 & 0\\
Fashion & 20 & Centroid & 8/8 & 0/0 & 99.316 & 92.664 & 1.120 & 0\\
Fashion & 20 & PCA & 8/8 & 0/0 & 99.219 & 87.671 & 1.366 & 0\\
MNIST & 3 & Centroid & 8/8 & 0/0 & 99.609 & 97.302 & 0.192 & 0\\
MNIST & 3 & PCA & 8/8 & 0/0 & 99.707 & 96.570 & 0.173 & 0\\
MNIST & 5 & Centroid & 32/32 & 0/0 & 99.707 & 96.613 & 0.397 & 0\\
MNIST & 5 & PCA & 32/32 & 0/0 & 99.609 & 95.990 & 0.369 & 0\\
MNIST & 10 & Centroid & 8/8 & 0/0 & 99.121 & 93.689 & 0.874 & 0\\
MNIST & 10 & PCA & 8/8 & 0/0 & 99.512 & 92.065 & 0.842 & 0\\
MNIST & 20 & Centroid & 8/8 & 0/0 & 98.926 & 89.294 & 1.410 & 1\\
MNIST & 20 & PCA & 8/8 & 0/0 & 98.633 & 85.730 & 1.463 & 3\\
\bottomrule
\end{tabular}
\end{adjustbox}
\end{table*}

\begin{table*}[tbp]
\centering
\caption{Native scaling on new MNIST seed-173 fits, using real held-out prefixes disjoint from 3,000 training records. Age 1 is the primary predeclared scale series; the additional age-8 prefixes evaluated by the implementation are retained as paired diagnostics. Both bases share the catalog and updates; prefixes are correlated. This measures numerical quality up to 16,384 records, not completed large-catalog proof speed.}
\label{tab:optimized-scale}
\begin{adjustbox}{max width=.98\textwidth,center=\textwidth}
\begin{tabular}{rrlrrrrrr}
\toprule
Age & $N$ & Basis & Decisions & Mean agr. & Worst agr. & Label cert. & Width & Exact rows\\
\midrule
1 & 128 & Centroid & 4/4 & 99.805 & 99.219 & 94.922 & 0.760 & 0\\
1 & 128 & PCA & 4/4 & 99.609 & 99.219 & 95.703 & 0.665 & 0\\
1 & 512 & Centroid & 4/4 & 99.951 & 99.805 & 95.703 & 0.742 & 0\\
1 & 512 & PCA & 4/4 & 99.805 & 99.609 & 95.996 & 0.649 & 0\\
1 & 1,024 & Centroid & 4/4 & 99.927 & 99.805 & 95.239 & 0.762 & 0\\
1 & 1,024 & PCA & 4/4 & 99.854 & 99.609 & 95.630 & 0.665 & 0\\
1 & 4,096 & Centroid & 4/4 & 99.847 & 99.780 & 95.123 & 0.730 & 0\\
1 & 4,096 & PCA & 4/4 & 99.854 & 99.805 & 95.496 & 0.639 & 0\\
1 & 16,384 & Centroid & 4/4 & 99.786 & 99.762 & 94.958 & 0.742 & 0\\
1 & 16,384 & PCA & 4/4 & 99.815 & 99.768 & 95.325 & 0.649 & 0\\
8 & 128 & Centroid & 4/4 & 100.000 & 100.000 & 96.875 & 0.287 & 0\\
8 & 128 & PCA & 4/4 & 99.609 & 99.219 & 95.703 & 0.306 & 0\\
8 & 512 & Centroid & 4/4 & 100.000 & 100.000 & 97.607 & 0.267 & 0\\
8 & 512 & PCA & 4/4 & 99.805 & 99.609 & 96.387 & 0.280 & 0\\
8 & 1,024 & Centroid & 4/4 & 100.000 & 100.000 & 97.095 & 0.291 & 0\\
8 & 1,024 & PCA & 4/4 & 99.805 & 99.609 & 95.801 & 0.306 & 0\\
8 & 4,096 & Centroid & 4/4 & 99.933 & 99.854 & 97.394 & 0.271 & 0\\
8 & 4,096 & PCA & 4/4 & 99.829 & 99.780 & 96.228 & 0.284 & 0\\
8 & 16,384 & Centroid & 4/4 & 99.904 & 99.860 & 97.446 & 0.267 & 0\\
8 & 16,384 & PCA & 4/4 & 99.850 & 99.829 & 96.336 & 0.278 & 0\\
\bottomrule
\end{tabular}
\end{adjustbox}
\end{table*}


The implementation computes exact residual energies using bounded signed
integer limbs, reuses center-only quadratic terms, and evaluates only each
row's authenticated group column in native switching bounds. A resumable
setup also obtains the proving-key size from its canonical saved bytes
instead of serializing the key a second time solely to count it. Hashing
streams large files. These changes preserve the proof relations, rounded
basis correction and integer outputs. New real-proof smoke checks verify
13 compact proofs and a separate 16-proof base with three refinement proofs;
all 20 associated transcript tampering checks reject. These are
implementation checks rather than an independent security audit.
Table~\ref{tab:optimized-kernels} measures native kernels only. Its ratios
are not substituted for end-to-end proof measurements.
\begin{table}[t]
\centering
\caption{Implementation-only paired native microbenchmarks, same 128-record catalog and exact integer outputs, three alternating-order repetitions. Times are median seconds. Ratios do not measure proof generation or deployment speed.}
\label{tab:optimized-kernels}
\begin{tabular}{lrrr}
\toprule
Kernel & Previous & Optimized & Ratio\\
\midrule
Residual statistics & 0.0862 & 0.0536 & 1.61\\
Center parameters & 0.4220 & 0.1471 & 2.87\\
Complete native bounds & 0.5584 & 0.2231 & 2.50\\
\bottomrule
\end{tabular}
\end{table}

